\documentclass[10pt,a4paper]{article}
\usepackage[a4paper,margin=14mm,top=12mm,bottom=14mm]{geometry}
\usepackage{fontspec}
\setmainfont{Liberation Sans}
\usepackage{polyglossia}
\setdefaultlanguage{french}
\usepackage{amsmath}
\usepackage{xcolor}
\usepackage{graphicx}
\usepackage[most]{tcolorbox}
\usepackage{tabularx,booktabs,array,multirow}
\usepackage{enumitem}
\usepackage{fancyhdr}
\usepackage{lastpage}
\usepackage[hidelinks]{hyperref}
\definecolor{uimm}{HTML}{005677}
\definecolor{cyan}{HTML}{13A2DD}
\definecolor{fond}{HTML}{E6F2FA}
\definecolor{ok}{HTML}{2E7D32}
\definecolor{ko}{HTML}{B3261E}
\definecolor{sylvain}{HTML}{7B1FA2}
\setlength{\parindent}{0pt}
\setlength{\parskip}{3pt}
\setlist{nosep}
\setlist[itemize]{leftmargin=13pt,label=\textcolor{cyan}{\textbullet}}
\newcommand{\titre}[1]{\vspace{5pt}{\color{uimm}\large\bfseries #1}\par\vspace{-3pt}{\color{cyan}\rule{\linewidth}{1.2pt}}\par\vspace{0pt}}
\tcbset{boite/.style={colback=fond,colframe=uimm,boxrule=0.6pt,arc=2pt,left=6pt,right=6pt,top=4pt,bottom=4pt},
        defi/.style={colback=sylvain!6,colframe=sylvain,boxrule=0.8pt,arc=2pt,left=6pt,right=6pt,top=4pt,bottom=4pt,
                     title={\bfseries Défi},coltitle=white,colbacktitle=sylvain,fonttitle=\small}}
\newcommand{\bandeau}[3]{%
\begin{tcolorbox}[colback=uimm,colframe=uimm,arc=2pt,left=4pt,right=8pt,top=4pt,bottom=4pt]
\begin{minipage}[c]{18mm}\setlength{\fboxsep}{1.2pt}\fcolorbox{white}{white}{\includegraphics[height=16.5mm]{logo_uimm.jpg}}\end{minipage}\hspace{4mm}%
\begin{minipage}[c]{\dimexpr\linewidth-22mm\relax}\color{white}
{\small #1}\\[1pt]{\LARGE\bfseries #2}\\[2pt]{\large #3}\hfill{\small CQPM Technicien qualité}
\end{minipage}
\end{tcolorbox}}
\newcommand{\reponse}[1][1]{\par\vspace{2pt}{\color{gray!60}\foreach \i in {1,...,#1}{\rule{\linewidth}{0.3pt}\par\vspace{9pt}}}}
\usepackage{pgffor}
\pagestyle{fancy}\fancyhf{}
\renewcommand{\headrulewidth}{0pt}
\newcommand{\pied}[1]{\fancyfoot[L]{\footnotesize\color{uimm}#1}\fancyfoot[R]{\footnotesize\color{uimm}Page \thepage{} sur \pageref{LastPage}}}

\renewcommand{\bandeau}[3]{%
\begin{tcolorbox}[colback=uimm,colframe=uimm,arc=2pt,left=4pt,right=8pt,top=4pt,bottom=4pt]
\begin{minipage}[c]{18mm}\setlength{\fboxsep}{1.2pt}\fcolorbox{white}{white}{\includegraphics[height=16.5mm]{logo_uimm.jpg}}\end{minipage}\hspace{4mm}%
\begin{minipage}[c]{\dimexpr\linewidth-22mm\relax}\color{white}
{\small #1}\\[1pt]{\LARGE\bfseries #2}\\[2pt]{\large #3}
\end{minipage}
\end{tcolorbox}}
\newcommand{\seance}[3]{%
\begin{tcolorbox}[colback=uimm,colframe=uimm,arc=2pt,left=6pt,right=6pt,top=3pt,bottom=3pt]
\color{white}{\Large\bfseries #1 · #2}\\[1pt]{\itshape #3}
\end{tcolorbox}}
\tcbset{pas/.style={colback=white,colframe=cyan,boxrule=0.6pt,arc=2pt,left=6pt,right=6pt,top=3pt,bottom=3pt,
                    title={\bfseries Dans Ellistat},coltitle=white,colbacktitle=cyan,fonttitle=\small},
        piege/.style={colback=ko!5,colframe=ko,boxrule=0.6pt,arc=2pt,left=6pt,right=6pt,top=3pt,bottom=3pt,
                    title={\bfseries Piège},coltitle=white,colbacktitle=ko,fonttitle=\small},
        verif/.style={colback=orange!6,colframe=orange!80!black,boxrule=0.6pt,arc=2pt,left=6pt,right=6pt,top=3pt,bottom=3pt,
                    title={\bfseries À vérifier sur poste avant la séance},coltitle=white,colbacktitle=orange!80!black,fonttitle=\small},
        corr/.style={colback=ok!5,colframe=ok,boxrule=0.6pt,arc=2pt,left=6pt,right=6pt,top=3pt,bottom=3pt,
                    title={\bfseries Corrigé},coltitle=white,colbacktitle=ok,fonttitle=\small}}
\newcommand{\fichier}[1]{\texttt{\small #1}}
\newcommand{\menu}[1]{\textsf{\textbf{#1}}}
\newcommand{\champ}[1]{\rule{#1}{0.3pt}}

\pied{Ellistat Data Analysis · six séances · livret stagiaire}
\begin{document}
\thispagestyle{fancy}
\bandeau{Travaux pratiques · Maîtrise statistique des procédés}{Six séances sur Ellistat}{Cartes de contrôle, capabilité et moyen de mesure avec Ellistat Data Analysis (version web)}

\begin{tcolorbox}[boite]
\textbf{La règle des six séances : prédire, faire, comparer.} Avant chaque calcul, vous écrivez ce que vous attendez. Ellistat calcule. Vous comparez avec un calcul à la main ou avec une page web de la formation. Quand les chiffres diffèrent, on cherche pourquoi : c'est là qu'on apprend. Aucune séance ne demande Excel : les fichiers sont prêts à importer.
\end{tcolorbox}

\titre{Le programme}
\renewcommand{\arraystretch}{1.25}
\begin{tabularx}{\linewidth}{@{}l X l@{}}
\toprule
\textbf{Séance} & \textbf{La question} & \textbf{Page}\\\midrule
S0 & Une grille, des lignes, des sous-groupes : comment Ellistat lit-il nos mesures ? & \pageref{s0}\\
S1 & La carte d'Ellistat dit-elle la même chose que la nôtre ? & \pageref{s1}\\
S2 & La carte sonne : que se passe-t-il, que fait-on ? & \pageref{s2}\\
S3 & On ne mesure pas, on compte : quelle carte, quelles limites ? & \pageref{s3}\\
S4 & Peut-on faire confiance à l'instrument et aux opérateurs ? & \pageref{s4}\\
S5 & Du moyen de mesure à la décision : un dossier complet & \pageref{s5}\\
\bottomrule
\end{tabularx}

\titre{Trois idées à garder}
\begin{itemize}
\item \textbf{Une ligne de la grille = un sous-groupe}, c'est-à-dire des pièces prélevées ensemble. Les pièces d'un sous-groupe sont rangées sur plusieurs colonnes côte à côte.
\item \textbf{« Répétition »} : réglage de la première de ces colonnes. Il indique combien de colonnes forment le sous-groupe. Le mot est trompeur : il a été pensé pour plusieurs mesures d'une même pièce, mais le calcul est le même pour 5 pièces consécutives.
\item \textbf{Court terme et global.} La dispersion à l'intérieur des sous-groupes donne Cp et Cpk. La dispersion de toutes les mesures donne Pp et Ppk. L'écart entre les deux révèle les dérives.
\end{itemize}

\titre{Les fichiers}
{\small Tous dans le dossier \fichier{donnees} de la page « Six séances sur Ellistat ». Format Excel, une feuille par fichier, importés directement dans Ellistat sans les ouvrir.}\par
\renewcommand{\arraystretch}{1.15}
\begin{tabularx}{\linewidth}{@{}>{\raggedright\arraybackslash}p{0.56\linewidth} X@{}}
\fichier{S0\_S1\_Butees\_6x5.xlsx} & 30 butées, 6 sous-groupes de 5 (S0, S1)\\
\fichier{S1\_Butees\_mauvais\_decoupage.xlsx} & les mêmes 30 butées, regroupées autrement (S1)\\
\fichier{S1\_S5\_Butees\_10x5.xlsx} & 50 butées, 10 sous-groupes de 5 (S1, S5)\\
\fichier{S2\_Pression.xlsx}, \fichier{S2\_Dimension.xlsx}, \fichier{S2\_Liquide.xlsx} & trois enquêtes (S2)\\
\fichier{S3\_Inspection.xlsx} & 120 lots de 150 pièces contrôlées (S3)\\
\fichier{S4\_Justesse\_Bosch\_p11.xlsx}, \fichier{S4\_RR\_Bosch\_p14.xlsx} & exemples du livret Bosch n° 10 (S4)\\
\fichier{S4\_Linearite\_simule.xlsx} & 5 cales, 10 mesures chacune, données simulées (S4)\\
\fichier{S5\_Justesse\_cale12\_simule.xlsx}, \fichier{S5\_RR\_butees\_simule.xlsx} & dossier de synthèse, données simulées (S5)\\
\end{tabularx}

% ------------------------------------------------------------------ S0
\newpage
\label{s0}
\seance{S0}{Prise en main}{Une grille, des lignes, des sous-groupes}

\textbf{Objectif.} Importer les butées dans un projet, régler la grille pour qu'Ellistat lise 6 sous-groupes de 5 pièces, et saisir la tolérance 11,85 / 12,00 mm.

\begin{tcolorbox}[pas]
\begin{enumerate}[label=\textbf{\arabic*.},leftmargin=14pt]
\item Page des projets : créer un projet à votre nom, par exemple « Butées Dupont ».
\item Ouvrir la grille du projet. Dans le menu en haut à droite de la grille : \menu{Importer depuis Excel}, puis choisir \fichier{S0\_S1\_Butees\_6x5.xlsx}.
\item Vérifier la grille : 6 lignes, 6 colonnes (Echantillon, P1 à P5). Chaque colonne porte un badge de type : les mesures doivent être numériques.
\item Ouvrir le panneau de la colonne \textbf{P1} : régler \menu{Répétition} sur 5, puis saisir tolérance min 11,85 et tolérance max 12,00.
\item \menu{Statistique descriptive} > \menu{Détail} : cocher la caractéristique. Ne pas cocher la colonne « Echantillon » : ce n'est qu'un numéro.
\end{enumerate}
\end{tcolorbox}

\textbf{Prédire avant d'ouvrir l'écran Détail}
\begin{enumerate}[label=\textcolor{cyan}{\textbf{Q\arabic*.}},leftmargin=22pt]
\item Combien de points doit compter la carte des moyennes ? Combien de pièces derrière chaque point ?\reponse
\item Dans cette grille, que représente une ligne ? Que représente une colonne ?\reponse
\end{enumerate}
\textbf{Constater}
\begin{enumerate}[label=\textcolor{cyan}{\textbf{Q\arabic*.}},leftmargin=22pt,resume]
\item La carte compte-t-elle le nombre de points prévu ? Relevez la moyenne générale affichée.\reponse
\item Remettez \menu{Répétition} à 1 sur P1. Que devient la carte ? Pourquoi ? Puis remettez 5.\reponse[2]
\item Le mot « Répétition » désigne d'habitude plusieurs mesures d'une même pièce. Pourquoi peut-on l'utiliser ici pour 5 pièces différentes ?\reponse[2]
\end{enumerate}

\begin{tcolorbox}[piege]
Un fichier qui contient des lignes vides (prévues pour des mesures futures) crée des sous-groupes vides dans la grille. Vérifiez toujours le nombre de lignes après l'import.
\end{tcolorbox}

% ------------------------------------------------------------------ S1
\newpage
\label{s1}
\seance{S1}{Comparer}{La carte d'Ellistat dit-elle la même chose que la nôtre ?}

\textbf{Objectif.} Calculer à la main les limites de la carte $\bar X / R$ des 30 butées, puis les comparer à celles d'Ellistat. Lire Cp, Cpk, Pp et Ppk. Comprendre pourquoi les chiffres diffèrent un peu, et pourquoi la décision ne change pas.

\textbf{A. À la main (calculatrice).} Complétez le tableau à partir de \fichier{S0\_S1\_Butees\_6x5.xlsx} (valeurs affichées dans la grille d'Ellistat).

\renewcommand{\arraystretch}{1.5}
\begin{tabularx}{\linewidth}{|l|*{6}{>{\centering\arraybackslash}X|}}\hline
\textbf{Sous-groupe} & 1 & 2 & 3 & 4 & 5 & 6\\\hline
Moyenne $\bar x$ & & & & & & \\\hline
Étendue $R$ & & & & & & \\\hline
\end{tabularx}

\vspace{3pt}
$\bar{\bar X}$ = \champ{20mm} \qquad $\bar R$ = \champ{20mm}\par\vspace{4pt}
$LCS_{\bar X} = \bar{\bar X} + 0{,}577\,\bar R$ = \champ{20mm} \quad $LCI_{\bar X}$ = \champ{20mm}\par\vspace{4pt}
$LCS_R = 2{,}114\,\bar R$ = \champ{18mm}

{\small Coefficients pour des sous-groupes de 5 pièces : $A_2 = 0{,}577$, $D_4 = 2{,}114$, $d_2 = 2{,}326$ (limites à 3 écarts-types).}

\textbf{B. Dans Ellistat.} Écran \menu{Détail} de la séance S0, référentiel \menu{Centré Moyenne}.

\renewcommand{\arraystretch}{1.4}
\begin{tabularx}{\linewidth}{|X|c|c|}\hline
 & \textbf{Votre calcul} & \textbf{Ellistat}\\\hline
Limite supérieure de la carte des moyennes & & \\\hline
Limite inférieure de la carte des moyennes & & \\\hline
Limite supérieure de la carte des étendues & & \\\hline
Écart-type court terme & $\bar R / d_2$ = & \\\hline
Cp = $(TS - TI) / (6\,\sigma_{\text{court terme}})$ & & \\\hline
Cpk & & \\\hline
Pp \quad / \quad Ppk & & \\\hline
\end{tabularx}

\begin{enumerate}[label=\textcolor{cyan}{\textbf{Q\arabic*.}},leftmargin=22pt]
\item Les limites diffèrent un peu. Les mêmes points sortent-ils des limites dans les deux calculs ?\reponse
\item Ellistat estime l'écart-type court terme autrement que $\bar R / d_2$. Lequel des deux utilise toutes les mesures du sous-groupe, et pas seulement la plus grande et la plus petite ?\reponse[2]
\item Cp = 1,22 mais Cpk est bien plus faible. Que dit l'histogramme sur le centrage ? Passez en \menu{Centré Cible} : que montre la carte ?\reponse[2]
\item Pp est plus petit que Cp. Que s'est-il passé entre les sous-groupes ?\reponse
\end{enumerate}

\textbf{C. Le mauvais découpage} (pour aller plus loin). Importez \fichier{S1\_Butees\_mauvais\_decoupage.xlsx} dans une nouvelle grille. Ce sont les mêmes 30 valeurs, mais chaque ligne regroupe la 1\textsuperscript{re} pièce (puis la 2\textsuperscript{e}…) de chacun des 6 prélèvements : des pièces fabriquées à des moments différents. Réglez \menu{Répétition} sur 6.
\begin{enumerate}[label=\textcolor{cyan}{\textbf{Q\arabic*.}},leftmargin=22pt,resume]
\item Prédiction : Cp va-t-il monter ou baisser ? Pp va-t-il changer ?\reponse
\item Relevez Cp et Pp. Expliquez avec la phrase : « la dérive est entrée dans le court terme ».\reponse[2]
\end{enumerate}

\newpage
\textbf{D. Cinquante butées.} Importez \fichier{S1\_S5\_Butees\_10x5.xlsx} (10 sous-groupes, Répétition = 5, même tolérance).
\begin{enumerate}[label=\textcolor{cyan}{\textbf{Q\arabic*.}},leftmargin=22pt,resume]
\item Quels sous-groupes sortent des limites ? Quelle forme a la courbe des moyennes ? Quelle cause peut-on imaginer au poste ?\reponse[2]
\item Comparez Cp, Cpk, Pp et Ppk. Lequel faut-il annoncer au client ? Pourquoi ne peut-on pas annoncer de capabilité tant que la carte n'est pas sous contrôle ?\reponse[2]
\item Sous les capabilités, Ellistat affiche Rs et Rr. Calculez Pp / Cp et Ppk / Pp. Que mesure chacun de ces deux rapports : la stabilité du procédé ou son réglage ?\reponse[2]
\end{enumerate}

\begin{tcolorbox}[boite]
\textbf{À retenir.} Cp dit ce que la machine sait faire à court terme. Pp dit ce qu'elle a fait sur toute la période. Ppk ajoute le centrage. Une capabilité ne s'annonce que sur un procédé sous contrôle.
\end{tcolorbox}

% ------------------------------------------------------------------ S2
\newpage
\label{s2}
\seance{S2}{Lire et décider}{La carte sonne : que se passe-t-il, que fait-on ?}

\textbf{Objectif.} En binôme, mener une enquête sur des données réelles, tracer la carte dans Ellistat, repérer les alertes, choisir une réaction, puis présenter au groupe en 5 minutes.

\textbf{E1 · La pression du réservoir} (\fichier{S2\_Pression.xlsx}). La pression d'un réservoir d'eau de refroidissement est relevée les 5 premières heures de chaque jour, du 1\textsuperscript{er} décembre au 13 janvier : 24 jours, une ligne par jour. Répétition = 5.

\textbf{E2 · Une dimension usinée} (\fichier{S2\_Dimension.xlsx}). Dimension critique d'une pièce usinée, en valeurs codées. 5 pièces consécutives toutes les heures, 27 sous-groupes. Répétition = 5.

\textbf{E3 · Un liquide fabriqué par lots} (\fichier{S2\_Liquide.xlsx}). Chaque lot est homogène : une seule valeur par lot, 17 lots. Ici, pas de sous-groupe : Répétition = 1. Ellistat trace une carte des valeurs individuelles.

\begin{tcolorbox}[pas]
Importer le fichier, régler la Répétition sur la première colonne de mesure, puis \menu{Statistique descriptive} > \menu{Détail}, référentiel \menu{Centré Moyenne}. Les points signalés par une règle portent le repère « Alerte » : tendance, même côté, proche des limites, mouvement alternatif.
\end{tcolorbox}

\begin{enumerate}[label=\textcolor{cyan}{\textbf{Q\arabic*.}},leftmargin=22pt]
\item Avant d'ouvrir la carte : dans la grille, quelles valeurs vous semblent anormales ?\reponse
\item Quels points sortent des limites ? Quels points portent une alerte, et laquelle ?\reponse[2]
\item Racontez ce qui s'est passé : un événement ponctuel, un changement durable, une dérive ?\reponse[2]
\item Les limites sont calculées sur toutes les données, y compris les périodes perturbées. Est-ce honnête ? Que faudrait-il faire pour surveiller la suite ?\reponse[2]
\item Quelle réaction proposez-vous au poste ? (ne rien toucher, plan de réaction, recentrer, prévoir une intervention, chercher ce qui a changé la dispersion)\reponse[2]
\end{enumerate}

\textbf{Présentation (5 minutes) :} la situation, la carte, ce qui sonne, votre décision.

% ------------------------------------------------------------------ S3
\newpage
\label{s3}
\seance{S3}{Compter les défauts}{On ne mesure pas, on compte : quelle carte, quelles limites ?}

\textbf{Objectif.} Choisir la carte aux attributs adaptée, la tracer dans Ellistat, vérifier ses limites à la main, et comprendre le rôle de la période de référence.

\textbf{La situation} (\fichier{S3\_Inspection.xlsx}). À la sortie d'une ligne, on contrôle des lots de 150 pièces. Pour chaque lot, on note le nombre de pièces non conformes : 120 lots.

\begin{enumerate}[label=\textcolor{cyan}{\textbf{Q\arabic*.}},leftmargin=22pt]
\item Compte-t-on des pièces défectueuses ou des défauts ? La taille des lots est-elle fixe ? Quelles cartes conviennent : p, np, c, u ?\reponse[2]
\end{enumerate}

\begin{tcolorbox}[pas]
\menu{Statistique descriptive} > \menu{Cartes de contrôle aux attributs}. Choisir la carte p, indiquer la taille de lot (150, ou la colonne Taille\_lot) et la colonne des non-conformes. Laisser d'abord la « Phase I (zone d'observation) » sur toutes les données.
\end{tcolorbox}

\begin{enumerate}[label=\textcolor{cyan}{\textbf{Q\arabic*.}},leftmargin=22pt,resume]
\item À la main : $\bar p = \dfrac{\text{total des non-conformes}}{\text{total des pièces}}$ = \champ{20mm}, puis $LCS = \bar p + 3\sqrt{\dfrac{\bar p(1-\bar p)}{n}}$ = \champ{20mm} avec $n = 150$. Comparez avec Ellistat.\reponse
\item Quels lots sortent ? Combien de non-conformes cela représente-t-il dans un lot ?\reponse
\item Tracez maintenant la carte np. Les mêmes lots sortent-ils ? Laquelle est la plus parlante pour un opérateur ?\reponse[2]
\item Réglez la Phase I sur les lots 1 à 20 seulement. Combien de lots sortent maintenant ? La période 1 à 20 était-elle représentative ?\reponse[2]
\end{enumerate}

\begin{tcolorbox}[defi]
\textbf{Carte P' de Laney.} Sur de gros lots, les limites classiques deviennent si serrées que trop de points sortent, alors que le procédé n'a pas changé : la variation entre lots dépasse ce que prévoit le modèle. Tracez la carte P' de Laney et comparez le nombre de lots hors limites. Laquelle des deux cartes croyez-vous ? Pourquoi ?
\end{tcolorbox}

{\small La page « Compter les défauts » de la formation aide à choisir la bonne carte (partie « Quelle carte pour quelle situation ? »).}

% ------------------------------------------------------------------ S4
\newpage
\label{s4}
\seance{S4}{Le moyen de mesure}{Peut-on faire confiance à l'instrument et aux opérateurs ?}

\textbf{Objectif.} Mener dans Ellistat les trois études du moyen de mesure : justesse (Cg, Cgk), linéarité, R\&R. Comparer avec les pages web de la formation, qui appliquent la méthode du livret Bosch n° 10.

\textbf{A. Justesse} (\fichier{S4\_Justesse\_Bosch\_p11.xlsx}). Une cale de valeur certifiée 6,002 mm, mesurée 50 fois. La cote contrôlée a une tolérance de 5,970 à 6,030 mm ; résolution 0,001 mm.
\begin{tcolorbox}[pas]
\menu{Gage R\&R} > \menu{Justesse} : valeur vraie $X_0$ = 6,002 ; tolérance min 5,970 et max 6,030 ; résolution 0,001 ; mesures prises dans la grille.
\end{tcolorbox}
\begin{enumerate}[label=\textcolor{cyan}{\textbf{Q\arabic*.}},leftmargin=22pt]
\item Relevez : écart de justesse \champ{18mm}, Cg \champ{14mm}, Cgk \champ{14mm}, tolérance minimale \champ{18mm}. Les voyants sont-ils verts ?\reponse
\item Collez les 50 mesures dans la page « L'instrument est-il capable ? » ($T = 0{,}060$, $x_m = 6{,}002$). Trouvez-vous les mêmes Cg et Cgk ? Sinon, quel facteur les sépare ?\reponse[2]
\item Cg est bon et Cgk plus faible : quel défaut l'instrument a-t-il ? Peut-on le corriger ?\reponse
\end{enumerate}

\textbf{B. Linéarité} (\fichier{S4\_Linearite\_simule.xlsx}, données simulées). Un micromètre 0-25 mm mesure 10 fois chacune de 5 cales : 5, 10, 15, 20 et 25 mm.
\begin{tcolorbox}[pas]
\menu{Gage R\&R} > \menu{Linéarité} : 5 produits, 10 valeurs par produit ; pour chaque cale, saisir sa valeur vraie puis ses 10 mesures (copier-coller depuis la page des séances).
\end{tcolorbox}
\begin{enumerate}[label=\textcolor{cyan}{\textbf{Q\arabic*.}},leftmargin=22pt,resume]
\item Prédiction : regardez les moyennes des 5 cales. Le biais est-il constant, ou change-t-il avec la longueur ?\reponse
\item Ellistat juge-t-il la pente significative ? Et la justesse ? Que faut-il faire de ce micromètre pour contrôler une cote de 12 mm ?\reponse[2]
\end{enumerate}

\textbf{C. R\&R} (\fichier{S4\_RR\_Bosch\_p14.xlsx}). 10 pièces, 3 opérateurs (A, B, C), 2 séries ; tolérance 5,970 à 6,030 mm.
\begin{tcolorbox}[pas]
\menu{Gage R\&R} > \menu{Mesures} > créer à partir de la grille. Une ligne par pièce, une colonne par couple opérateur et série (A\_r1 … C\_r2). Ne pas prendre la colonne « Piece » comme mesure. Méthode \textbf{ANOVA}.
\end{tcolorbox}
\begin{enumerate}[label=\textcolor{cyan}{\textbf{Q\arabic*.}},leftmargin=22pt,resume]
\item Relevez : répétabilité \champ{16mm}, reproductibilité \champ{16mm}, interaction \champ{16mm}, \%GRR \champ{14mm}, ndc \champ{10mm}.\reponse
\item Le livret Bosch annonce \%GRR = 17,95 \%, comme la page « L'opérateur compte-t-il ? » (bouton « Exemple du livret Bosch »). Trouvez-vous la même valeur ? Sinon, d'où vient l'écart ?\reponse[2]
\item Le système de mesure est-il capable, capable sous conditions ou non capable ? Qui pèse le plus : l'instrument ou les opérateurs ?\reponse[2]
\end{enumerate}

% ------------------------------------------------------------------ S5
\newpage
\label{s5}
\seance{S5}{Synthèse}{Du moyen de mesure à la décision : un dossier complet}

\textbf{La situation.} Une butée de diamètre 11,85 / 12,00 mm. Le responsable qualité vous confie trois fichiers et une question : « Peut-on piloter ce poste avec une carte de contrôle, et le procédé est-il capable ? » Les données des études de mesure sont simulées.

\renewcommand{\arraystretch}{1.25}
\begin{tabularx}{\linewidth}{@{}l l X@{}}
\toprule
\textbf{Étape} & \textbf{Fichier} & \textbf{Ce qu'il contient}\\\midrule
1. Justesse & \fichier{S5\_Justesse\_cale12\_simule.xlsx} & 50 mesures d'une cale de 12,000 mm composée de deux cales (10 + 2), micromètre au 0,001 mm\\
2. R\&R & \fichier{S5\_RR\_butees\_simule.xlsx} & 10 butées, 3 opérateurs, 2 séries\\
3. Production & \fichier{S1\_S5\_Butees\_10x5.xlsx} & 50 butées, 10 sous-groupes de 5\\
\bottomrule
\end{tabularx}

\textbf{La consigne.} Un projet à votre nom avec trois grilles. Respectez l'ordre : on ne juge pas une machine avec un instrument qu'on n'a pas validé. Partagez le projet avec le formateur (\menu{utilisateurs autorisés}).

\renewcommand{\arraystretch}{1.5}
\begin{tabularx}{\linewidth}{|l|X|X|}\hline
\textbf{Étape} & \textbf{Résultats clés} & \textbf{Conclusion}\\\hline
1. Justesse & Cg = \hspace{14mm} Cgk = & \\\hline
2. R\&R & \%GRR = \hspace{9mm} ndc = & \\\hline
3a. Carte & points hors limites : & \\\hline
3b. Capabilité & Cp = \hspace{6mm} Cpk = \hspace{6mm} Pp = \hspace{6mm} Ppk = & \\\hline
\end{tabularx}

\textbf{Votre décision} (5 lignes, à présenter au groupe) : le moyen de mesure, l'état du procédé, la capabilité que l'on peut annoncer, l'action à mener en premier.\reponse[5]

\begin{tcolorbox}[piege]
Les données d'entreprise que vous mettriez dans Ellistat partent sur un serveur extérieur. Ne le faites qu'avec l'accord écrit de votre entreprise, dans un projet personnel ni public ni partagé.
\end{tcolorbox}

\vfill
{\footnotesize\color{uimm}\textbf{Sources.} Ellistat, guide utilisateur de Data Analysis (pages consultées le 30/10/2026) et réponse d'Ellistat du 10/11/2026 sur les sous-groupes. Robert Bosch GmbH, \emph{Booklet 10 : Capability of Measurement and Test Processes}, édition novembre 2019, exemples p. 11 et p. 14. Données S2 et S3 : jeux SPC-2, SPC-3, SPC-10 et SPC-1 de la formation. Données simulées signalées comme telles.}
\end{document}
